\documentclass[11pt]{article}
\usepackage[T1]{fontenc}
\usepackage[margin=1in]{geometry}
\usepackage{amsmath,amssymb,mathtools}
\usepackage[colorlinks=true,linkcolor=blue,citecolor=blue,urlcolor=blue]{hyperref}

\title{Noether charge from the finite perturbative action}
\author{}
\date{}

\newcommand{\dd}{\mathrm{d}}
\newcommand{\cL}{\mathcal{L}}
\newcommand{\Lg}{\mathcal{L}_{\mathrm g}}
\newcommand{\cH}{\mathcal{H}}
\newcommand{\cE}{\mathcal{E}}
\newcommand{\cR}{\mathcal{R}}
\newcommand{\Lie}{\mathcal{L}}

\begin{document}
\maketitle

This fragment fills in the Noether-charge calculation for the finite
perturbative action
\begin{align}
  I^{(0)}[h,k,\varphi]
  &=\frac12\int_M\dd^3x\,\sqrt{-g^{(0)}}
    \left(\Lg^{(1)}[k]+\Lg^{(2)}[h]\right)
    \notag\\
  &\quad
    -\frac12\int_M\dd^3x\,\sqrt{-g^{(0)}}
    \left(\nabla_\mu\varphi\nabla^\mu\varphi+m^2\varphi^2\right).
\end{align}
All indices are raised and lowered with the AdS$_3$ background
$g^{(0)}_{\mu\nu}$, and $\nabla^{(0)}_\mu$ is the corresponding covariant
derivative. For this local calculation $\xi^\mu$ is taken to be a background
Killing vector. The extension to asymptotic Killing vectors requires the
separate boundary falloff analysis.

The first variation is written as
\begin{align}
  \delta I^{(0)}
  &=\int_M\dd^3x\,\sqrt{-g^{(0)}}
    \left(
      E^{(h)}_{\mu\nu}\delta h^{\mu\nu}
      +E^{(k)}_{\mu\nu}\delta k^{\mu\nu}
      +E^{(\varphi)}\delta\varphi
    \right) \notag\\
  &\quad
    +\int_{\Sigma_f-\Sigma_i}
    \dd^2x\,\sqrt{\sigma^{(0)}}\,\tau_\mu\Theta^\mu,
\end{align}
where
\begin{align}
  E^{(h)}_{\mu\nu}
  &=-\frac12\left(
    -h_{\mu\nu}
    -\frac12\nabla^{(0)2}h_{\mu\nu}
    +\frac12\nabla^{(0)}_\mu\nabla^{(0)}_\rho h_\nu{}^\rho
    +\frac12\nabla^{(0)}_\nu\nabla^{(0)}_\rho h_\mu{}^\rho
    \right. \notag\\
  &\qquad\left.
    -\frac12\nabla^{(0)}_\mu\nabla^{(0)}_\nu h
    +\frac12g^{(0)}_{\mu\nu}\nabla^{(0)2}h
    -\frac12g^{(0)}_{\mu\nu}
      \nabla^{(0)}_\rho\nabla^{(0)}_\sigma h^{\rho\sigma}
    \right),\\
  E^{(k)}_{\mu\nu}&=0,\\
  E^{(\varphi)}&=\nabla^2\varphi-m^2\varphi.
\end{align}
The symplectic potential current is
\begin{align}
  \Theta^\mu
  &=-\frac12\left(
    \nabla^{(0)}_\nu\delta k^{\mu\nu}
    -\nabla^{(0)\mu}\delta k
    \right)
    +\nabla^\mu\varphi\,\delta\varphi \notag\\
  &\quad
    -\frac12\left(
    h^{\mu\nu}\nabla^{(0)}_\nu\delta h
    -\frac12h\nabla^{(0)\mu}\delta h
    +h^{\nu\rho}\nabla^{(0)\mu}\delta h_{\nu\rho}
    -h^{\nu\rho}\nabla^{(0)}_\rho\delta h^\mu{}_\nu
    \right. \notag\\
  &\qquad
    -h^{\mu\nu}\nabla^{(0)}_\rho\delta h_\nu{}^\rho
    +\frac12h\nabla^{(0)}_\nu\delta h^{\mu\nu}
    +\frac12\nabla^{(0)\mu}h_{\nu\rho}\delta h^{\nu\rho}
    \notag\\
  &\qquad\left.
    +\frac12\nabla^{(0)}_\nu h\,\delta h^{\mu\nu}
    -\nabla^{(0)}_\rho h^{\nu\rho}\delta h^\mu{}_\nu
    \right).
\end{align}

The field-space vector generated by $\xi^\mu$ is
\begin{align}
  X_\xi
  =\int\dd^3x\left(
    \Lie_\xi h_{\mu\nu}\frac{\delta}{\delta h_{\mu\nu}}
    +\Lie_\xi k_{\mu\nu}\frac{\delta}{\delta k_{\mu\nu}}
    +\Lie_\xi\varphi\frac{\delta}{\delta\varphi}
  \right).
\end{align}
Acting on the action gives
\begin{align}
  X_\xi\cdot\delta I^{(0)}
  =-\int_{\Sigma_f-\Sigma_i}
    \dd^2x\,\sqrt{\sigma^{(0)}}\,\tau_\mu\xi^\mu
    \cL[h,k,\varphi].
\end{align}
Thus the Noether charge is
\begin{align}
  H_\xi
  &=\int_\Sigma\dd^2x\,\sqrt{\sigma^{(0)}}\,\tau_\mu
    \left(X_\xi\cdot\Theta^\mu+\xi^\mu\cL[h,k,\varphi]\right) \notag\\
  &=H_{\xi,h}+H_{\xi,k}+H_{\xi,\varphi}.
\end{align}

\section*{The $h_{\mu\nu}$ contribution}

The quadratic metric contribution is
\begin{align}
  H_{\xi,h}
  =\int_\Sigma\dd^2x\,\sqrt{\sigma^{(0)}}\,
    \tau_\mu\cH^\mu_{\xi,h},
\end{align}
with
\begin{align}
  \cH^\mu_{\xi,h}
  &=-\frac12\,\mathcal{A}^\mu_{\xi,h}
    +\frac12\xi^\mu\Lg^{(2)}[h],
\end{align}
where
\begin{align}
  \mathcal{A}^\mu_{\xi,h}
  &=h^{\mu\nu}\nabla^{(0)}_\nu\Lie_\xi h
    -\frac12h\nabla^{(0)\mu}\Lie_\xi h
    +h^{\nu\rho}\nabla^{(0)\mu}\Lie_\xi h_{\nu\rho}
    -h^{\nu\rho}\nabla^{(0)}_\rho\Lie_\xi h^\mu{}_\nu
    \notag\\
  &\quad
    -h^{\mu\nu}\nabla^{(0)}_\rho\Lie_\xi h_\nu{}^\rho
    +\frac12h\nabla^{(0)}_\nu\Lie_\xi h^{\mu\nu}
    +\frac12\nabla^{(0)\mu}h_{\nu\rho}\Lie_\xi h^{\nu\rho}
    \notag\\
  &\quad
    +\frac12\nabla^{(0)}_\nu h\,\Lie_\xi h^{\mu\nu}
    -\nabla^{(0)}_\rho h^{\nu\rho}\Lie_\xi h^\mu{}_\nu .
\end{align}
This current can be decomposed as
\begin{align}
  \cH^\mu_{\xi,h}
  =\xi_\nu T_{(h)}^{\mu\nu}
    +\nabla^{(0)}_\nu S_\xi^{\mu\nu}
    +\cR^\mu_{\xi,h}.
\end{align}
The local gravitational-energy representative is
\begin{align}
  T_{(h)}^{\mu\nu}=-\cE^{(2)\mu\nu}[h,h],
\end{align}
where
\begin{align}
  \cE_{\mu\nu}^{(2)}[h,h]
  =R_{\mu\nu}^{(2)}[h,h]
   -\frac12g^{(0)}_{\mu\nu}R^{(2)}[h,h]
   -\frac12h_{\mu\nu}R^{(1)}[h].
\end{align}
For any symmetric tensor $X_{\mu\nu}$,
\begin{align}
  R_{\mu\nu}^{(2)}[X,X]
  &=\nabla^{(0)}_\rho\Gamma^{(2)}[X,X]^\rho{}_{\nu\mu}
    -\nabla^{(0)}_\nu\Gamma^{(2)}[X,X]^\rho{}_{\rho\mu}
    \notag\\
  &\quad
    +\Gamma^{(1)}[X]^\rho{}_{\rho\lambda}
      \Gamma^{(1)}[X]^\lambda{}_{\nu\mu}
    -\Gamma^{(1)}[X]^\rho{}_{\nu\lambda}
      \Gamma^{(1)}[X]^\lambda{}_{\rho\mu},\\
  R^{(2)}[X,X]
  &=g^{(0)\mu\nu}R_{\mu\nu}^{(2)}[X,X]
    -X^{\mu\nu}R_{\mu\nu}^{(1)}[X]
    -2X_{\mu\nu}X^{\mu\nu},\\
  R_{\mu\nu}^{(1)}[X]
  &=\frac12\left(
    \nabla^{(0)}_\rho\nabla^{(0)}_\mu X^\rho{}_\nu
    +\nabla^{(0)}_\rho\nabla^{(0)}_\nu X^\rho{}_\mu
    -\nabla^{(0)2}X_{\mu\nu}
    -\nabla^{(0)}_\mu\nabla^{(0)}_\nu X
    \right),\\
  R^{(1)}[X]
  &=\nabla^{(0)}_\mu\nabla^{(0)}_\nu X^{\mu\nu}
    -\nabla^{(0)2}X+2X,\\
  \Gamma^{(2)}[X,X]^\rho{}_{\mu\nu}
  &=-\frac12X^{\rho\sigma}
    \left(
      \nabla^{(0)}_\mu X_{\sigma\nu}
      +\nabla^{(0)}_\nu X_{\mu\sigma}
      -\nabla^{(0)}_\sigma X_{\mu\nu}
    \right),\\
  \Gamma^{(1)}[X]^\rho{}_{\mu\nu}
  &=\frac12\left(
    \nabla^{(0)}_\mu X^\rho{}_\nu
    +\nabla^{(0)}_\nu X^\rho{}_\mu
    -\nabla^{(0)\rho}X_{\mu\nu}
  \right).
\end{align}
The antisymmetric surface current is
\begin{align}
  S_\xi^{\mu\nu}
  &=s_2\nabla^{(0)[\mu}\xi^{\nu]}
    +\frac12hD_{1,\xi}^{[\mu\nu]}
    +D_{2,\xi}^{[\mu\nu]},\\
  s_2
  &=-\frac14h_{\rho\sigma}h^{\rho\sigma}+\frac18h^2,\\
  D_{1,\xi}^{\mu\nu}
  &=-h^{\mu\rho}\nabla^{(0)}_\rho\xi^\nu
    +\Gamma^{(1)}[h]^{\nu\mu}{}_\rho\xi^\rho,\\
  D_{2,\xi}^{\mu\nu}
  &=h^\mu{}_\lambda h^{\lambda\rho}
    \nabla^{(0)}_\rho\xi^\nu
    -h^{\mu\rho}\Gamma^{(1)}[h]^\nu{}_{\rho\sigma}\xi^\sigma
    +\Gamma^{(2)}[h,h]^{\nu\mu}{}_\rho\xi^\rho.
\end{align}
The remainder is proportional to the linearized Einstein equation:
\begin{align}
  \cR^\mu_{\xi,h}
  =\xi^\rho h^{\mu\sigma}\cE^{(1)}_{\rho\sigma}[h]
   -\frac12h\xi_\nu\cE^{(1)\mu\nu}[h],
\end{align}
with
\begin{align}
  \cE_{\mu\nu}^{(1)}[X]
  &=-X_{\mu\nu}
    -\frac12\nabla^{(0)2}X_{\mu\nu}
    +\frac12\nabla^{(0)}_\mu\nabla^{(0)}_\rho X_\nu{}^\rho
    +\frac12\nabla^{(0)}_\nu\nabla^{(0)}_\rho X_\mu{}^\rho
    \notag\\
  &\quad
    -\frac12\nabla^{(0)}_\mu\nabla^{(0)}_\nu X
    +\frac12g^{(0)}_{\mu\nu}\nabla^{(0)2}X
    -\frac12g^{(0)}_{\mu\nu}
      \nabla^{(0)}_\rho\nabla^{(0)}_\sigma X^{\rho\sigma}.
\end{align}
Thus on linearized solutions, $\cE^{(1)}_{\mu\nu}[h]=0$, the bulk part of
$H_{\xi,h}$ is represented by $\xi_\nu T_{(h)}^{\mu\nu}$ up to the surface
term $\nabla^{(0)}_\nu S_\xi^{\mu\nu}$.

\section*{The $k_{\mu\nu}$ contribution}

The second-order backreaction field contributes
\begin{align}
  H_{\xi,k}
  =\int_\Sigma\dd^2x\,\sqrt{\sigma^{(0)}}\,\tau_\mu
    \cH^\mu_{\xi,k},
\end{align}
where
\begin{align}
  \cH^\mu_{\xi,k}
  &=-\frac12\left(
      \nabla^{(0)}_\nu\Delta_\xi k^{\mu\nu}
      -\nabla^{(0)\mu}\Delta_\xi k
    \right)
    +\frac12\xi^\mu\Lg^{(1)}[k] \notag\\
  &=\nabla^{(0)}_\nu S_{\xi,k,0}^{\mu\nu},
\end{align}
with
\begin{align}
  S_{\xi,k,0}^{\mu\nu}
  =\xi^{[\mu}
    \left(
      \nabla^{(0)}_\rho k^{\nu]\rho}
      -\nabla^{(0)\nu]}k
    \right).
\end{align}
Thus the $k_{\mu\nu}$ sector enters this calculation as a surface term.

\section*{The scalar contribution}

The scalar field gives the ordinary matter charge density
\begin{align}
  H_{\xi,\varphi}
  &=\int_\Sigma\dd^2x\,\sqrt{\sigma^{(0)}}\,\tau_\mu
    \cH^\mu_{\xi,\varphi},\\
  \cH^\mu_{\xi,\varphi}
  &=\xi_\nu T_{(\varphi)}^{\mu\nu},
\end{align}
where
\begin{align}
  T_{(\varphi)}^{\mu\nu}
  =\nabla^{(0)\mu}\varphi\nabla^{(0)\nu}\varphi
   -\frac12g^{(0)\mu\nu}
    \left(
      \nabla^{(0)}_\rho\varphi\nabla^{(0)\rho}\varphi
      +m^2\varphi^2
    \right).
\end{align}

Combining the three pieces gives
\begin{align}
  H_\xi
  &=\int_\Sigma\dd^2x\,\sqrt{\sigma^{(0)}}\,\tau_\mu
    \left[
      \xi_\nu T_{(h)}^{\mu\nu}
      +\xi_\nu T_{(\varphi)}^{\mu\nu}
      +\nabla^{(0)}_\nu
       \left(S_\xi^{\mu\nu}+S_{\xi,k,0}^{\mu\nu}\right)
      +\cR^\mu_{\xi,h}
    \right].
\end{align}
On solutions of the first-order graviton equation,
\begin{align}
  \cE^{(1)}_{\mu\nu}[h]=0,
\end{align}
the remainder $\cR^\mu_{\xi,h}$ vanishes. The second-order Einstein equation
is
\begin{align}
  \cE^{(1)}_{\mu\nu}[k]
  +\cE^{(2)}_{\mu\nu}[h,h]
  =T_{(\varphi)\mu\nu}.
\end{align}
Since $T_{(h)}^{\mu\nu}=-\cE^{(2)\mu\nu}[h,h]$, this constraint is equivalently
\begin{align}
  \cE^{(1)\mu\nu}[k]
  =T_{(h)}^{\mu\nu}+T_{(\varphi)}^{\mu\nu}.
\end{align}
Therefore, on the first-order equation and the second-order constraint,
\begin{align}
  H_\xi
  =\int_\Sigma\dd^2x\,\sqrt{\sigma^{(0)}}\,\tau_\mu
    \left[
      \xi_\nu\cE^{(1)\mu\nu}[k]
      +\nabla^{(0)}_\nu
       \left(S_\xi^{\mu\nu}+S_{\xi,k,0}^{\mu\nu}\right)
    \right].
\end{align}
For a background Killing vector in unit AdS$_3$,
\begin{align}
  \xi_\nu\cE^{(1)\mu\nu}[k]
  =\nabla^{(0)}_\nu\widehat S_{\xi,k}^{\mu\nu},
\end{align}
where antisymmetrization brackets have unit weight and
\begin{align}
  \widehat S_{\xi,k}^{\mu\nu}
  &=\xi^{[\mu}\nabla^{(0)\nu]}k
    -\xi^{[\mu}\nabla^{(0)}_\rho k^{\nu]\rho}
    +\xi_\rho\nabla^{(0)[\mu}k^{\nu]\rho}
    +\frac12 k\nabla^{(0)[\mu}\xi^{\nu]}
    -k^{\rho[\mu}\nabla^{(0)}_\rho\xi^{\nu]}.
\end{align}
Defining the total $k$-sector surface current by
\begin{align}
  S_{\xi,k}^{\mu\nu}
  &:=S_{\xi,k,0}^{\mu\nu}+\widehat S_{\xi,k}^{\mu\nu} \notag\\
  &=\xi_\rho\nabla^{(0)[\mu}k^{\nu]\rho}
    +\frac12 k\nabla^{(0)[\mu}\xi^{\nu]}
    -k^{\rho[\mu}\nabla^{(0)}_\rho\xi^{\nu]},
\end{align}
the total charge is represented purely by a boundary surface current:
\begin{align}
  H_\xi
  =\int_\Sigma\dd^2x\,\sqrt{\sigma^{(0)}}\,\tau_\mu
    \nabla^{(0)}_\nu
    \left(S_\xi^{\mu\nu}+S_{\xi,k}^{\mu\nu}\right).
\end{align}
Equivalently, with $n_\mu$ the outward unit normal to $\partial\Sigma$ inside
$\Sigma$ and
$\dd\Sigma^{(0)}_{\mu\nu}
  :=\dd x\,\sqrt{\gamma^{(0)}}\,\tau_{[\mu}n_{\nu]}$,
\begin{align}
  H_\xi
  =\int_{\partial\Sigma}\dd\Sigma^{(0)}_{\mu\nu}
    \left(S_\xi^{\mu\nu}+S_{\xi,k}^{\mu\nu}\right).
\end{align}

\end{document}
